\documentclass[10pt]{article}
\usepackage{amsmath,amsfonts,amssymb,mathtools,mathrsfs,bm,bbm,array,xcolor,geometry}
\allowdisplaybreaks
\setlength{\parindent}{0pt}
\setlength{\parskip}{4pt}
\emergencystretch=4em

\geometry{paperwidth=215mm,paperheight=290mm,hmargin=15mm,vmargin=16mm}
\begin{document}
\section*{P0154 Train timetables and speed profiles for commuter and metro services}
\subsection*{Mathematical formulation}
\[
\min\sum_{k\in K}t_{k,s_k^e}.
\]

\[
t_{k,\widehat s}=t_{ks}+n_{ks}.
\]

\[
t_{is}\ge t_{js}+n_{js}+\delta-\lambda_{ij}^sM,\qquad
t_{js}\ge t_{is}+n_{is}+\delta-(1-\lambda_{ij}^s)M.
\]

\begin{align*}
t_{is}&\ge t_{js}+h^{min}-\lambda_{ij}^sM,&
t_{js}&\ge t_{is}+h^{min}-(1-\lambda_{ij}^s)M,\\
t_{is}&\le t_{js}+h^{max}+\lambda_{ij}^sM+(1-x_{ji})M,&
t_{js}&\le t_{is}+h^{max}+(1-\lambda_{ij}^s)M+(1-x_{ij})M.
\end{align*}

\[
\sum_{j\ne i}x_{ij}\le1,\quad \sum_{i\ne j}x_{ij}\le1,\quad
\sum_i\sum_{j\ne i}x_{ij}=|K_s|-1,\quad x_{ij}\le\lambda_{ij}^s.
\]

\[
v^{in}_{k,s_k^b}=v^{in}_{k,s_k^e}=0,\quad
0\le v^{in}_{ks},v^{cru}_{ks}\le v_s^{lim},\quad
t_{ks}\le t^{bc}_{ks}\le t^{ec}_{ks}\le t_{k,\widehat s}.
\]

\begin{align*}
v^{cru}_{ks}&=v^{in}_{ks}+a_{ks}(t^{bc}_{ks}-t_{ks}),\\
v^{in}_{k,\widehat s}&=v^{cru}_{ks}+b_{ks}(t_{k,\widehat s}-t^{ec}_{ks}),\\
L_s&=\tfrac12(v^{in}_{ks}+v^{cru}_{ks})(t^{bc}_{ks}-t_{ks})
 +v^{cru}_{ks}(t^{ec}_{ks}-t^{bc}_{ks})
 +\tfrac12(v^{cru}_{ks}+v^{in}_{k,\widehat s})(t_{k,\widehat s}-t^{ec}_{ks}).
\end{align*}

\[
\check\theta=\mathbf1_{\{t_{i,\widehat{s'}}\ge t_{js}\}},\quad
\hat\theta=\mathbf1_{\{t_{i,\widehat{s'}}\le t_{j,\widehat s}\}},\quad
\zeta=\check\theta\hat\theta\lambda_{ij}^{s'}.
\]

\begin{align*}
t_{js}-t_{i,\widehat{s'}}&\le M(1-\check\theta),&
t_{js}-t_{i,\widehat{s'}}&\ge1-(M+1)\check\theta,\\
t_{i,\widehat{s'}}-t_{j,\widehat s}&\le M(1-\hat\theta),&
t_{i,\widehat{s'}}-t_{j,\widehat s}&\ge1-(M+1)\hat\theta,\\
\zeta&\ge\check\theta+\hat\theta+\lambda_{ij}^{s'}-2,&
\zeta&\le\check\theta,\ \hat\theta,\ \lambda_{ij}^{s'} .
\end{align*}

\[
t_{j,\widehat s}-t_{i,\widehat{s'}}
\ge \frac{v^{max}_{js}}{b^{max}_{emergency}}+\delta
 -(3-\iota_{\alpha,js}-\iota_{\beta,js}-\zeta)M.
\]

\begin{align*}
\iota_{1,ks}=1&\ \Longleftrightarrow\ v^-\le v^0,&
\iota_{2,ks}=1&\ \Longleftrightarrow\ v^-=v^0,&
\iota_{3,ks}=1&\ \Longleftrightarrow\ v^-\ge v^0,\\
\iota_{4,ks}=1&\ \Longleftrightarrow\ v^0\le v^+,&
\iota_{5,ks}=1&\ \Longleftrightarrow\ v^0=v^+,&
\iota_{6,ks}=1&\ \Longleftrightarrow\ v^0\ge v^+ .
\end{align*}

\[
\mathcal I=\{(1,4),(3,6),(2,5),(1,5),(2,6),(1,6)\}.
\]

\[
v^{max}_{js}=\max\{v^{in}_{js},v^{cru}_{js},v^{in}_{j,\widehat s}\}.
\]

\subsection*{Definitions and model scope}
A railway operator schedules ordinary and through trains $K$ in a commuter and metro network with some shared tracks. Each train's path, section running times, section lengths, speed limits, minimum and maximum departure headways, safety allowance, normal acceleration and braking capabilities and emergency braking capability are given. Tracks assigned to independent directions do not interfere. The operator chooses section-entry, acceleration-end and braking-start times $t_{ks}$, $t^{bc}_{ks}$ and $t^{ec}_{ks}$ as nonnegative integer seconds, along with nonnegative entry and cruising speeds and constant acceleration and braking acceleration for each stage. Acceleration lies between zero and its given maximum, while braking acceleration lies between the negative normal braking capability and zero. The aim is the smallest sum of all trains' terminal arrival times.

Entry to a successor section equals entry to the current section plus fixed running time. Every ordered pair of trains sharing a physical section has a zero-or-one precedence record $\lambda_{ij}^s$, with the reverse record equal to one minus the forward record. On each nonmoving-block section, for each pair in given identifier order, the first train's entry is at least the second's entry plus its running time and the safety allowance, less the large time bound multiplied by the precedence record. Conversely, the second's entry is at least the first's entry plus its running time and the allowance, less that bound times one minus the record. All times lie in a supplied finite range sufficient for full operation, and the large bound also covers running and safety intervals.

For trains sharing an origin, each departure must exceed the other's departure by at least the minimum headway, allowing respective reductions of the large time bound times the precedence record and times one minus that record. Direct-successor records $x_{ij}$ take only zero or one, with one exactly when the second train departs immediately after the first from their common origin and zero otherwise; these records apply only to distinct trains sharing that origin. Each train's outgoing successor-record sum and incoming predecessor-record sum are at most one. All such records at an origin sum to its train count less one, and each cannot exceed the corresponding precedence record. The first departure less the second cannot exceed the maximum headway plus the large bound times precedence plus the large bound times one minus the reverse direct-successor record. The opposite upper limit uses one minus precedence and one minus the forward direct-successor record.

For each train, sum the direct-successor records multiplied by whether each successor is a through train. Add, over every two-link chain whose second successor is neither the train itself nor its first successor, the product of both successor records and the second successor's through-service flag. This gives the subsequent-through count $B_i$. For an ordinary train it must equal its outgoing direct-successor-record sum; for a through train it cannot exceed one. At the sequence end, only successors that actually exist are counted.

Every train has zero entry speed at its origin and terminal boundary, and all entry and cruising speeds obey the section limit. Acceleration ends no earlier than section entry, braking begins no earlier than acceleration ends and no later than section exit. Cruising speed equals entry speed plus section acceleration times acceleration duration. Entry speed on the next section equals cruising speed plus braking acceleration times braking duration. Section length equals the average of the acceleration-stage endpoint speeds times its duration, plus cruising speed times cruising duration, plus the average of braking-stage endpoint speeds times braking duration. Accelerations may differ between sections; intermediate boundaries need not be stops, and stage durations are not fixed at twenty seconds.

Safety is checked only for the supplied shared-track combinations $\mathcal Q$. Fixed-block checks use adjacent fixed-block sections, moving-block checks use the same moving-block section, and independent tracks do not create check combinations. For each combination, the first time flag is one exactly when the leading train leaves its specified checking section no earlier than the following train enters the checked section. The second flag is one exactly when that leading exit is no later than the following train's exit. The overlap-and-order flag $\zeta$ is one exactly when both time flags and the relevant precedence record are all one, and otherwise zero. With integer seconds, a zero time flag requires the reversed comparison to differ by at least one second, while equality belongs to the one case.

Each speed profile has six zero-or-one comparison records $\iota_{1,ks}$, $\iota_{2,ks}$, $\iota_{3,ks}$, $\iota_{4,ks}$, $\iota_{5,ks}$ and $\iota_{6,ks}$. They are one exactly when entry speed is no greater than, equal to or no less than cruising speed, and when cruising speed is no greater than, equal to or no less than exit speed, respectively. At equality, the applicable weak-comparison and equality records are simultaneously one. Every checking combination separately uses the record pairs first and fourth, third and sixth, second and fifth, first and fifth, second and sixth, and first and sixth. The following train's exit less the leading train's checking-section exit must be at least the maximum of the following train's three speeds divided by emergency braking capability plus the safety allowance, less the large time bound multiplied by three minus the two speed records and the overlap-and-order flag. The maximum is calculated from the actual profile. These requirements retain finite allowances when a pair is inactive. Requiring an unreduced safety interval for every overlap is only an optional conservative operating choice.

\subsection*{Inline notation and relations}
$K$; $S_k$; $\mathcal F$; $Z$; $S_B$; $\widehat s$; $s$; $n_{ks}$; $\xi_k=1$; $t_{ks},t^{bc}_{ks},t^{ec}_{ks}$; $v^{in}_{ks},v^{cru}_{ks}\ge0$; $a_{ks}\in[0,a^{max}]$; $b_{ks}\in[-b^{max}_{service},0]$; $\lambda_{ij}^s\in\{0,1\}$; $\lambda_{ji}^s=1-\lambda_{ij}^s$; $x_{ij}\in\{0,1\}$; $j$; $i$; $i,j\in K_s,\ i\ne j,\ s\in S_B$; $s\notin Z$; $i<j$; $s=s_i^b=s_j^b$; $K_s$; $B_i=\sum_{j\ne i}x_{ij}\xi_j+
\sum_{j\ne i}\sum_{j'\notin\{i,j\}}x_{ij}x_{jj'}\xi_{j'}$; $B_i=\sum_{j\ne i}x_{ij}$; $B_i\le1$; $\mathcal Q$; $(i,j,s',s)$; $s'=\widehat s$; $s'=s\in Z$; $\epsilon=1$; $v^{max}_{js}$; $\iota_\alpha,\iota_\beta$; $\zeta=1$; $M$; $\sigma$; $\sigma=0$; $k,s$; $v^-=v^{in}_{ks}$; $v^0=v^{cru}_{ks}$; $v^+=v^{in}_{k,\widehat s}$; $(\alpha,\beta)\in\mathcal I$; $v^{in}_{j,\widehat s},v^{in}_{js},v^{in}_{js},
v^{in}_{j,\widehat s},v^{in}_{js},v^{cru}_{js}$; $s'$; $a\ge0,b\le0$; $(1,6)$; $0\le t\le\sum_sn_s+(|K|+1)h^{max}$.
\end{document}